\input{../article_base.tex}
\title{פתרון מטלה 03 --- משוואות דיפרנציאליות, 80320}
% chktex-file 9
% chktex-file 17

\DeclareMathOperator{\sgn}{sgn}

\begin{document}
\maketitle
\maketitleprint[blue]{}

\question{}
יהי $0 \in I \subseteq \RR$ קטע פתוח ו־$A : I \to M_n(\RR), g : I \to \RR^n$ רציפות.
נראה כי לכל $x_0$ קיים פתרון יחיד $x : I \to \RR^n$ המוגדר על כל $I$ לבעיית ההתחלה,
\[
	\begin{cases}
		x' = A x + g \\
		x(0) = x_0
	\end{cases}
\]
\begin{proof}
	נבחין כי זוהי משוואה לא אוטונומית המוגדרת על־ידי $F(x, t) = A x + t$, בתרגיל הקודם ראינו כי קיים פתרון יחיד בסביבה של $0$ למשוואה אם היא ליפשיץ מקומית בציר הראשון וליפשיץ באותו קבוע בכל הציר השני.
	קיימת סביבה סגורה של $0$ ובה נקבל מהרחבה של משפט קנטור שהפונקציה רציפה במידה שווה בתחום ולכן בפרט $1$־ליפשיץ ובהתאם ממשפט פיקארד נובע שאכן קיימת סביבה פתוחה של $0$ בה יש פתרון יחיד לבעיה.

	עתה נניח ש־$x : J_{x_0}^* \to \RR^n$ פתרון מקסימלי של המשוואה, אז תנאי משפט הטריכוטומיה של זרימות חלים.
\end{proof}

\question{}
יהי $\alpha \in \RR$ ונגדיר את הבעיה,
\[
	\begin{cases}
		x' = \alpha x - x^3 \\
		x(0) = x_0
	\end{cases}
	\tag{1}
\]

\subquestion{}
נראה שאם $\alpha > 0$ אז קיימים בדיוק שלושה פתרונות קבועים ל־$(1)$ ושאם $\alpha \le 0$ קיים פתרון קבוע יחיד לבעיה.
\begin{proof}
	נבחין כי פתרון $x$ הוא קבוע אם ורק אם $x' = 0$, כלומר אם $x (\alpha - x^2) = 0$, בהתאם נקבל ש־$x = 0$ הוא פתרון קבוע ללא תלות ב־$\alpha$.
	אם $\alpha \ge 0$ אז נקבל שגם $\alpha - x^2 = 0$ מקבל שני פתרונות, $x = \pm \sqrt{\alpha}$, ולכן מצאנו שיש לפחות שלושה פתרונות ופתרון יחיד לבעיה בהתאמה.

	מפתרונות המשוואה נוכל להסיק שאם $\alpha \ge 0, x \notin \{0, \pm \sqrt{\alpha} \}$ אז $x (\alpha - x^2) \ne 0$ ולכן בפרט גם $x' \ne 0$ והפתרון המתקבל לא יהיה קבוע.
	אם $\alpha < 0$ ו־$x \ne 0$ אז נקבל שוב $x \ne 0, \alpha - x^2 < 0$ ולכן גם הפעם הפתרון לא יהיה קבוע.
\end{proof}

\subquestion{}
נוכיח שכל פתרון של $(1)$ הוא פונקציה מונוטונית.
\begin{proof}
	עבור פונקציות קבועות הטענה ברורה.
	נבחין כי המשוואה נתונה על־ידי פונקציה $F(x) = \alpha x - x^3$ היא גזירה ברציפות (ואף חלקה),
	ולכן ממשפט הטריכוטומיה וממשפט פיקארד נקבל שאם פתרון הוא לא קבוע אז בהכרח $x(t) \notin \{ 0, \pm \sqrt{\alpha} \}$ או $x(t) \ne 0$ לכל $t$, שכן אחרת היינו מקבלים התלכדות עם פונקציה קבועה.

	נניח ש־$\alpha \ge 0$ וכן $x(t) > \sqrt{\alpha} \iff x^2 - \alpha > 0$, אז נקבל $x' = x (\alpha - x^2) < 0$, כלומר $x$ מונוטונית יורדת ממש.
	אם $0 < x < \sqrt{\alpha}$ אז נקבל $x' = x (\alpha - x^2) > 0$ ולכן הפונקציה עולה ממש.
	עבור $-\sqrt{\alpha} < x < 0$ ו־$x < -\sqrt{\alpha}$ נקבל מהלכים כמעט זהים וש־$x$ מונוטונית יורדת ועולה ממש בהתאמה.

	אם $\alpha < 0$ ו־$x > 0$ אז נקבל $x > 0, \alpha - x^2 < 0$ ולכן גם $x' < 0$ ומונוטוניות יורדת ממש.
	אם $x < 0$ אז נקבל $x < 0, \alpha - x^2 < 0$ ולכן $x' > 0$, כלומר הפעם קיבלנו מונוטוניות עולה ממש.
\end{proof}

\subquestion{}
נמצא את ערכי $x_0, \alpha$ עבורם לפתרון של $(1)$ יש נקודת פיתול.
\begin{solution}
	ניזכר כי ל־$x$ יש נקודת פיתול ב־$t$ אז $x''(t) = 0$.
	נחשב ונקבל,
	\[
		x''(t) = 0
		\iff \alpha x'(t) - 3 x^2(t) x'(t) = 0
	\]
	לפונקציות קבועות אין נקודת פיתול ולכן נוכל להסיק שיכולה להיות נקודת פיתול בנקודות $\alpha = 3 x^2(t) \iff x(t) = \pm \sqrt{\frac{\alpha}{3}}$.

	מצאנו שהפונקציה תמיד מונוטונית וחסומה ולכן יש לה גבול ומתקיים, ולכן גם לנגזרת יש גבול כתוצאה של קיום המשוואה, אבל אם פונקציה מתכנסת וגם נגזרתה מתכנסת, אז הנגזרת מתכנסת לאפס, כלומר,
	\[
		0
		= \lim_{t \to \infty} x'(t)
		= \lim_{t \to \infty} \alpha x(t) - x^3(t)
		= \alpha L - L^3
	\]
	כאשר $L = \lim_{t \to \infty} x(t)$. אז $L (1 - L^2) = 0$.
	נקבל שלכל בחירת $\alpha, x_0$ יש רק ערך $L$ אחד שיתכן, לדוגמה עבור $\alpha > 0, x > \sqrt{\alpha}$ נקבל ש־$L = \sqrt{\alpha}$ בלבד בשל החסימה מלמטה והמונוטוניות היורדת.
	בהתאם עבור $\alpha \ge 0, x_0 > \sqrt{\alpha}$ נקבל שתמיד $x > \sqrt{\frac{\alpha}{3}}$ ולכן אין לפונקציה נקודת פיתול כלל.
	באופן דומה נקבל זאת עבור $\alpha \ge 0, x_0 < -\sqrt{\alpha}$, אבל מערך ביניים נקבל שכן יש נקודת פיתול עבור $\alpha \ge 0, 0 < x_0 \sqrt{\alpha}$ ועבור $-\sqrt{\alpha} < x < 0$.
	עבור $\alpha < 0$ נקבל ש־$\lim_{t \to \infty} x(t) = 0$ תמיד ולכן אם $|x_0| \ge \sqrt{\frac{\alpha}{3}}$ נקבל שנקודת פיתול היא מובטחת.
	לבסוף עבור שאר המקרים נקבל ש־$|x'(t)| = |\alpha x(t) - x^3(t)| > |x(t)| < 0$ ומשיקולים שראינו עבור המשוואה $x = x'$ נסיק ש־$\lim_{x \to -\infty} |x(t)| = \infty$ ולכן מערך ביניים יש נקודת פיתול בהכרח.
\end{solution}

\question{}
יהי $t_0 > 0$ ו־$a > 0$ ונגדיר את המשוואה,
\[
	\begin{cases}
		\frac{x'}{1 + x^2} + \frac{x}{t} = 1 \\
		x(t_0) = a
	\end{cases}
	\tag{1}
.\]

\subquestion{}
נסביר למה קיים פתרון ל־$(1)$ בסביבה של $t_0$.
\begin{proof}
	זוהי משוואה אוטונומית וגזירה ברציפות בסביבה של $t_0$ המוגדרת על־ידי המשוואה $F(x, t) = (1 + x^2)(1 - \frac{x}{t})$, מגזירות ברציפות נסיק שהיא ליפשיץ מקומית ולכן יש לה ממשפט פיקארד פתרון יחיד בסביבה של $t_0$.
\end{proof}

\subquestion{}
נסמן את תחום ההגדרה של הפתרון $x$ של $(1)$ ב־$J = (T_-, T_+)$.
נראה שלכל $t \in (t_0, T_+)$ מתקיים $x(t) > 0$.
\begin{proof}
	נניח שב־$t \in (t_0, T_+)$ מסוים מתקבל $x(t) \le 0$.
	בהתאם מערך הממוצע קיימת נקודה $t_0 < s < t$ עבורה מתקיים $0 > x'(s) = F(x(s), s)$,
	אז,
	\[
		0 > (1 + x^2(s)) (1 - \frac{x(s)}{s})
		\iff 0 > 1 - \frac{x(s)}{s}
		\iff x(s) > s
	.\]
	בחרנו את $s$ להיות נקודה בה הנגזרת שלילית, עתה נבחין כי חייבת להיות נקודה כזו כך ש־$x(s) < 0$, ישירות מסופיות מספר הנקודות שעוברות ב־$0$ וערך ממוצע, נניח ש־$s$ שבחרנו מקיימת זאת.
	נובע אם כך ש־$x(s) \le 0$ ולכן אבל $x(s) > s > 0$ בסתירה.
\end{proof}

\subquestion{}
נוכיח כי קיים $T_0 \in J$ כך שמתקיים $x(T_0) \le T_0$.
\begin{proof}
	בסעיף הקודם מצאנו ש־$x(t) > 0$ אם ורק אם לא קיים $x(s) > s$, כלומר קיימת לפחות נקודה אחת המקיימת $x(s) \ge s$, נסמן אותה ב־$T_0$.
\end{proof}

\subquestion{}
נראה שלכל $t \in (T_0, T_+)$ מתקיים $x(t) \le t$ ונסיק ש־$T_+ = \infty$.
\begin{proof}
	מתקיים,
	\[
		x(t) \le t
		\iff 0 \le 1 - \frac{x(t)}{t}
		\iff 0 \le F(x(t), t)
		\iff 0 \le x'(t)
	.\]
	נניח בשלילה שקיימת נקודה $s$ כך ש־$x(s) > s$, מהשקילות שראינו זה עתה $x'(s) < 0$.
	בהתאם נקבל שהפונקציה $1 - \frac{x(t)}{t}$ עולה בסביבה של $s$.
	נבחר $s_0 = \inf\{ t \in (T_0, T_+) \mid x'(t) < 0 \}$, אז מתקיים $x'(s_0) = 0$ שכן אחרת $s_0 = T_0$ בסתירה.
	בהתאם ישנה סביבה של $s_0$ בה $x$ מונוטונית עולה בסתירה להגדרתה.

	קיבלנו שאכן $x(t) \le t$ ב־$(T_0, T_+)$.
	ממשפט הטריכוטומיה נסיק שאם $T_+ < \infty$ אז $\limsup_{t \to T_+^-} x(t) \le T_+$ בסתירה לתנאי $\lVert x(t) \rVert \to \infty$, ולכן $T_+ = \infty$ בלבד.
\end{proof}

\subquestion{}
נראה שמתקיים,
\[
	\limsup_{t \to \infty} \frac{x(t)}{t} = 1
.\]
\begin{proof}
	בסעיף הקודם מצאנו ש־$x'(t) \ge 0$ לכל $t > T_0$, כלומר $x$ מונוטונית עולה החל מ־$T_0$.
	גם מצאנו שמתקיים $0 \le x(t) \le t$ ולכן $0 \le \frac{x(t)}{t} \le 1$, ונניח בשלילה שהגבול הוא לא $1$, לכן קיים $\varepsilon > 0$ כך שהחל ממקום מסוים $\frac{x(t)}{t} < 1 - \varepsilon$.
	נקבל ש־$\varepsilon < 1 - \frac{x(t)}{t}$, ולכן באופן דומה לסעיפים הקודמים,
	\[
		x'(t) > \varepsilon
		\iff F(x(t), t) > \varepsilon
		\iff (1 + x^2(t)) (1 - \frac{x(t)}{t}) > \varepsilon
	.\]
	אבל $(1 + x^2) \ge 1$ ולכן קיבלנו את הרצוי, דהינו $x'(t) > \varepsilon$, ולכן גם $\frac{x(t)}{t} > \varepsilon$ מאינטגרציה.
	נסיק,
	\[
		x'(t)
		= (1 + x^2(t)) (1 - \frac{x(t)}{t})
		> (1 + x^2(t)) \varepsilon
		\ge (1 + \varepsilon^2 t^2) \varepsilon
	.\]
	ונקבל ש־$x'(t) \to \infty$ בסתירה, לכן $\varepsilon = 0$ בלבד, כלומר הגבול של $\frac{x(t)}{t}$ הוא $1$ בלבד.
\end{proof}

\question{}
תהי $F : \RR^n \to \RR^n$ ליפשיץ מקומית ונניח שקיים $C > 0$ כך שלכל $x \in \RR^n$ מתקיים,
\[
	\lVert F(x) \rVert
	\le C(\lVert x \rVert + 1)
.\]
נגדיר את המשוואה,
\[
	\begin{cases}
		x' = F \circ x \\
		x(0) = x_0
	\end{cases}
	\tag{1}
.\]
נראה כי הפתרון למשוואה מוגדר לכל $\RR$.
\begin{proof}
	אין לתחום שפה ולכן ממשפט הטריכוטומיה אם הפתרון לא מוגדר ב־$\RR$ אז $\lVert x(t) \rVert \xrightarrow{t \to T_+} \infty$ כאשר $x : (T_-, T_+) \to \RR^n$.
\end{proof}

\end{document}
